\begin{textsample}{POS Dim 3 – df – Score 10.00 – df/df\_inf\_23.txt}  \label{ex:f3_pos_108}
16.2.
Corpo, \textbf{gestos} e movimentos Autor : Arthur Di Fábio Castro Esse campo de experiência propõe o trabalho voltado ao desenvolvimento corporal da \textbf{criança} que, ao se expressar, interage com o mundo desde cedo por meio de \textbf{gestos} e movimentos corporais, sejam eles dotados de intencionalidade ou de impulsos próprios da infância, bem como de espontaneidade ou coordenação de movimentos, \textbf{gestos} e sentidos.
A \textbf{criança} \textbf{brinca} e interage em diversas situações sociais e culturais as quais está exposta, estabelecendo relações que produzem conhecimentos sobre si e o outro e, \textbf{progressivamente}, tomando consciência de sua corporeidade.
Na Educação \textbf{Infantil}, as linguagens se entrelaçam e as diversas dimensões de aprendizagem se fundem na expressão da \textbf{criança}, o que torna essencial o trabalho corporal como instrumento de interação e comunicação que possibilita seu desenvolvimento e aprendizagem.
O trabalho corporal educativo na Educação \textbf{Infantil} deve levar em conta a centralidade do corpo da \textbf{criança}, voltando -o para o conhecimento e reconhecimento de suas potencialidades, limites, sensações e funções corporais.
Dessa forma, o corpo, como veículo de expressão das diversas linguagens ( a música, a dança, o teatro e as \textbf{brincadeiras}, dentre outras ), comunica -se com outros campos de experiência, de modo a promover possibilidades de desenvolvimento integral.
Nesse processo, é fundamental considerar ainda as contribuições de todas as matrizes culturais que compõem a sociedade brasileira.
Assim, jogos e \textbf{brincadeiras} de origem africana, indígena e europeia, que deram origem à população brasileira, por exemplo, devem ser considerados para o planejamento das ações na Educação \textbf{Infantil}.
Os \textbf{cuidados} físicos necessários com o corpo perpassam as interações da \textbf{criança} com o meio, com o outro e consigo mesma, fato que torna o trabalho educativo corporal primordial ao desenvolvimento da noção do que é seguro ou do que pode promover riscos para sua integridade física.
No entanto, ressalta -se que tais \textbf{cuidados} devem propiciar à \textbf{criança} condições de expressão sem que supostas limitações tolham seu desenvolvimento.
O trabalho pedagógico nesse campo de experiência deve propiciar explorações de movimentos que envolvam o próprio repertório da \textbf{criança}, ampliando -o à descoberta de variados modos de ocupação dos espaços com o corpo, bem como de atividades que lhe possibilite expressões cognitivas e afetivas em suas relações sociais e culturais, entrelaçadas às diversas linguagens e campos de experiências trabalhados.
Para tal, o repertório deve abranger atividades que envolvam \textbf{mímica}, expressões faciais e gestuais ; sonoridades ; olhares ; sentar com apoio ; rastejar, engatinhar, escorregar e caminhar, \textbf{apoiando} -se ou \textbf{livremente} ; correr ; alongar ; escalar ; saltar ; dar cambalhotas ; equilibrar -se e rolar.
Além dessas, o repertório pode incluir também as atividades que surgirem das \textbf{brincadeiras} e interações propostas no trabalho educativo com outras linguagens e campos de experiência, em que a autonomia e o protagonismo \textbf{infantil} devem ser levados em consideração nos objetivos pretendidos nesse campo de experiência.

% matched lemmas: apoiar, brincadeira, brincar, criança, cuidado, gesto, infantil, livremente, mímico, progressivamente
\end{textsample}
